This blueprint accompanies the Lean formalization of
\emph{On monic abelian cubic polynomials}, by Shubhrajit Bhattacharya and
Andrew O'Desky. It follows our November 2026 revision, which extends the
trace-one results of \href{https://arxiv.org/abs/2310.17831}{the arXiv version}
to arbitrary trace.

The statements below are taken from the revised paper. Each has links to its
Lean statement and proof, followed by a short explanation of the argument.
The numbering here is local to the blueprint; the TeX labels retained in the
source identify the corresponding results in the paper.

Our proof uses Poisson summation and the critical-line
mean-square estimate of Chandrasekharan and Narasimhan. Both are explicit
hypotheses in Lean. The comparison with Santens's theorem and the alternative
ideal count use the additional inputs listed in the last chapter.

For the strip estimate, Lean uses the residue theorem in place of
Titchmarsh's result. For the coefficient sum, it uses Perron's formula and
contour shifting in place of Alberts's theorem. We describe these arguments
at the corresponding places. For the moduli theorem, we formalize the
rational-point correspondence used in counting; the assertion over general
bases is not formalized. .

The \href{dep_graph_document.html}{dependency graph} is extracted from
the Lean proofs. It records references between the final registered declarations.
Some proofs use shared helper lemmas directly, so an isolated node need not
represent mathematics unused elsewhere. Green nodes have formalized statements and proofs,
relative to any literature inputs shown. Nodes for assumed inputs have
formalized statements and no proofs. The
\href{https://github.com/aodesky/abelian-cubics/blob/main/CORRESPONDENCE.md}{correspondence record}
gives further details.
